problem:  
Let \( G \) be a compact connected Lie group, \( H \subset G \) a closed subgroup such that \( M = G/H \) is a compact simply connected homogeneous space. Denote by \( \mathcal{M}_G^+(M) \) the space of \(G\)-invariant Riemannian metrics of **strictly positive sectional curvature**, and by \( \mathcal{M}_G^{\ge 0}(M) \) those of **nonnegative sectional curvature**.  

Each metric \( g \in \mathcal{M}_G^{\ge 0}(M) \) can be seen as a point in a finite-dimensional cone determined by the decomposition of \(\mathfrak{g}/\mathfrak{h}\) into irreducible \(\operatorname{Ad}(H)\)-modules.  

**Question:**  
For which homogeneous spaces \(M = G/H\) is the closure of \( \mathcal{M}_G^{+}(M) \) in this cone equal to all of \( \mathcal{M}_G^{\ge 0}(M) \)?  

Equivalently, does every \(G\)-invariant nonnegatively curved homogeneous metric occur as a limit (in the parameter space of invariant metrics) of positively curved ones—possibly through rescalings on submodules of \(\mathfrak{g}/\mathfrak{h}\)?  

If not universally true, classify those pairs \((G,H)\) where the equality fails.  

Why is it a "good" problem:  
This question investigates the **geometric topology of the curvature map** in the space of homogeneous metrics, probing the continuous border between strictly positive and merely nonnegative curvature in a fully invariant setup. It is a “good’’ problem for several deep reasons:

1. **Conceptual simplicity and depth:**  
   The statement is simple—asking whether the subset of positive-curvature homogeneous metrics is dense among nonnegative ones—but resolving it would require deep understanding of invariant curvature operators and their degenerations.

2. **Bridges differential geometry and Lie theory:**  
   Determining the closure structure of \( \mathcal{M}_G^{+}(M) \) demands combining analytic curvature inequalities with algebraic data from the isotropy representation decomposition, intertwining Riemannian geometry, invariant theory, and representation theory in a nontrivial way.

3. **Extends classical classification in a subtle direction:**  
   The classical positive curvature classifications (Berger, Wallach, Bérard-Bergery) identify where \( \mathcal{M}_G^{+}(M) \neq \emptyset \). This question instead asks about the *boundary behavior* of these metrics and whether limits preserve curvature sign—a fundamentally new layer of the classification program.

4. **Potentially transformative consequences:**  
   - If the closure coincidence holds universally, it implies a kind of **“metric stability” of positive curvature** under invariant degenerations.
   - If it fails, discovering explicit counterexamples would reveal new geometric strata within the moduli of invariant metrics, likely with subtle geometric transitions (e.g., emergence of flat directions in curvature operator spectra).

Thus the problem is both natural and profound: its resolution would reorganize the understanding of how curvature sign varies within invariant families, linking algebraic structure and geometric stability—a quintessential hallmark of a “good’’ research problem.